% --- SECCIÓ 3: OBJECTIUS I HIPÒTESI ---
\apartat{Objectius i Hipòtesi}

\subsection{Objectius del Projecte}
El present treball es proposa assolir una sèrie de fites tant teòriques com pràctiques:
\begin{itemize}
    \item \textbf{Objectiu 1:} Analitzar les bases científiques de la comunicació no verbal i la seva importància en el desenvolupament humà.
    \item \textbf{Objectiu 2:} Estudiar el funcionament dels sistemes de Comunicació Augmentativa i Alternativa (CAA) actuals.
    \item \textbf{Objectiu 3:} Dissenyar i construir un prototip funcional de taula sensorial que integri estímuls tàctils, visuals i sonors.
    \item \textbf{Objectiu 4:} Avaluar la utilitat de la taula com a eina de suport per a col·lectius amb necessitats especials.
\end{itemize}

\newpage
\subsection{Hipòtesi de Recerca}
La hipòtesi principal que guia aquest treball és la següent:
\begin{quote}
    \textit{"L'ús d'una taula sensorial dissenyada específicament per a l'estimulació multisensorial facilita la intenció comunicativa i millora la capacitat d'expressió de les persones amb dificultats greus en la parla."}
\end{quote}

Aquesta hipòtesi se sustenta en la teoria de la integració sensorial, que defensa que un cervell ben regulat sensorialment és més capaç de processar informació social complexa.
